% ============================================================
%  C: Como Programar -- 6ª Edição | Deitel & Deitel
%  Capítulo 10 -- Estruturas, Uniões, Manipulações de Bits
%               e Enumerações em C
%  FAZENDO A DIFERENÇA (10.21)
% ============================================================
\documentclass[12pt,a4paper]{article}
\usepackage[utf8]{inputenc}
\usepackage[T1]{fontenc}
\usepackage[brazil]{babel}
\usepackage{geometry}
\geometry{top=2.5cm,bottom=2.5cm,left=3cm,right=3cm}
\usepackage{parskip}\setlength{\parskip}{6pt}\setlength{\parindent}{0pt}
\usepackage{xcolor,tcolorbox,enumitem,listings,fancyhdr,titlesec,hyperref,booktabs,array,amsmath}
\tcbuselibrary{skins,breakable}

\definecolor{deitelblue}{RGB}{0,70,127}
\definecolor{deitelgray}{RGB}{240,242,245}
\definecolor{deitelgreen}{RGB}{0,110,60}
\definecolor{deitellightblue}{RGB}{220,235,250}
\definecolor{deitelred}{RGB}{160,20,20}
\definecolor{deitellorange}{RGB}{200,90,0}
\definecolor{ecogreen}{RGB}{0,120,50}
\definecolor{ecolightgreen}{RGB}{210,240,220}

\newtcolorbox{boaprática}[1][]{colback=deitellightblue,colframe=deitelblue,
  fonttitle=\bfseries\small,title={Boa prática de programação},breakable,#1}
\newtcolorbox{respostabox}[1][]{colback=deitelgray,colframe=deitelblue!60,
  fonttitle=\bfseries\small\color{deitelblue},title={Resposta detalhada},breakable,#1}
\newtcolorbox{enunciadoBox}[1][]{colback=white,colframe=deitelblue,
  fonttitle=\bfseries,title={#1},breakable}
\newtcolorbox{saida}[1][]{colback=black!88,colframe=deitelblue,
  fontupper=\ttfamily\small\color{green!70},
  title={\small\color{white}Saída do programa},breakable,#1}
\newtcolorbox{pesquisa}[1][]{colback=ecolightgreen,colframe=ecogreen,
  fonttitle=\bfseries\small,title={Pesquisa externa realizada},breakable,#1}
\newtcolorbox{atencao}[1][]{colback=yellow!10,colframe=deitellorange,
  fonttitle=\bfseries\small,title={Atenção},breakable,#1}

\setlist[enumerate,1]{label=\textbf{\alph*)},leftmargin=2em}
\setlist[itemize]{leftmargin=2em}

\lstset{
  basicstyle=\ttfamily\small,backgroundcolor=\color{deitelgray},
  frame=single,framerule=0.4pt,rulecolor=\color{deitelblue!40},
  breaklines=true,showstringspaces=false,
  keywordstyle=\color{deitelblue}\bfseries,
  commentstyle=\color{deitelgreen}\itshape,
  stringstyle=\color{deitelred},
  language=C,numbers=left,numberstyle=\tiny\color{gray},
  xleftmargin=18pt,xrightmargin=5pt,stepnumber=1,
}

\pagestyle{fancy}\fancyhf{}
\fancyhead[L]{\small\color{deitelblue}\textbf{C: Como Programar -- 6ª ed. | Deitel \& Deitel}}
\fancyhead[R]{\small\color{deitelblue}Cap. 10 -- Fazendo a Diferença}
\fancyfoot[C]{\small\thepage}
\renewcommand{\headrulewidth}{0.4pt}

\titleformat{\section}[block]{\large\bfseries\color{deitelblue}}{\thesection.}{0.5em}{}[\titlerule]
\titleformat{\subsection}[block]{\normalsize\bfseries\color{deitelblue!80}}{}{}{}

\hypersetup{colorlinks=true,linkcolor=deitelblue,urlcolor=ecogreen}

\begin{document}

\begin{titlepage}
  \centering\vspace*{2cm}
  {\color{deitelblue}\rule{\linewidth}{2pt}}\\[0.4cm]
  {\huge\bfseries\color{deitelblue}C: Como Programar}\\[0.2cm]
  {\large\color{deitelblue}6ª Edição $\cdot$ Paul Deitel \& Harvey Deitel}\\[0.2cm]
  {\color{deitelblue}\rule{\linewidth}{2pt}}\\[1cm]
  {\LARGE\bfseries Capítulo 10}\\[0.3cm]
  {\Large\color{deitelblue!80}Estruturas, Uniões, Manipulações de Bits e Enumerações}\\[0.8cm]
  \begin{tcolorbox}[colback=ecolightgreen,colframe=ecogreen,width=0.85\linewidth]
    \centering{\large\bfseries\color{ecogreen}Fazendo a Diferença}\\[0.3cm]
    {\normalsize Exercício 10.21\\[0.2cm]
    Informatização de registros de saúde\\
    Sistema integrado com IMC e frequência cardíaca}
  \end{tcolorbox}
  \vfill{\small Abordagem incremental Deitel}
\end{titlepage}

\tableofcontents\newpage

% ============================================================
\section{Exercício 10.21 -- Perfil de Saúde}
% ============================================================

\begin{enunciadoBox}[Enunciado 10.21]
Crie uma estrutura \texttt{PerfilSaude} com nome, sobrenome, sexo, data de
nascimento, altura (cm) e peso (kg). Implemente funções que:
\begin{itemize}[noitemsep]
  \item Recebam esses dados e definam os membros da estrutura.
  \item Calculem a idade em anos.
  \item Calculem a frequência cardíaca máxima e ideal.
  \item Calculem o IMC.
\end{itemize}
Exiba todas as informações junto com a tabela de classificação de IMC.
\end{enunciadoBox}

\subsection*{Pesquisa Externa}

\begin{pesquisa}
\textbf{Frequência cardíaca máxima:} fórmula clássica de Tanaka
\[
\text{FC}_{\max} = 220 - \text{idade}
\]

\textbf{Frequência cardíaca ideal (zona aeróbica):} 50\% a 85\% da FC\textsubscript{máx}.
A faixa típica para exercícios de queima de gordura é 60--70\%; para
condicionamento cardiovascular, 70--85\%.

\textbf{IMC:} $\text{IMC} = \text{peso (kg)} / \text{altura (m)}^2$, com classificação
da OMS (abaixo do peso, normal, sobrepeso, obeso) -- a mesma do Ex.~2.32.

\textbf{Cálculo de idade:} diferença entre o ano atual e o ano de nascimento,
\textit{ajustada} se o aniversário ainda não ocorreu no ano corrente.
\end{pesquisa}

\subsection*{Solução em C}

\begin{respostabox}
\begin{lstlisting}
/* Ex10.21: perfil_saude.c
   Sistema integrado de perfil de saude */
#include <stdio.h>
#include <string.h>
#include <time.h>

/* Estrutura para data */
typedef struct {
    int dia;
    int mes;
    int ano;
} Data;

/* Estrutura principal de perfil */
typedef struct {
    char nome[ 30 ];
    char sobrenome[ 30 ];
    char sexo;            /* 'M' ou 'F' */
    Data nascimento;
    double altura_cm;
    double peso_kg;
} PerfilSaude;

/* Prototipos */
void   lerPerfil       ( PerfilSaude *p );
int    calcularIdade   ( const Data *nasc );
int    fcMaxima        ( int idade );
double fcIdealMin      ( int idade );
double fcIdealMax      ( int idade );
double calcularIMC     ( double peso_kg, double altura_cm );
const char* classificarIMC( double imc );
void   exibirPerfil    ( const PerfilSaude *p );

int main( void )
{
    PerfilSaude paciente;

    printf( "===== SISTEMA DE PERFIL DE SAUDE =====\n\n" );
    lerPerfil( &paciente );
    exibirPerfil( &paciente );

    return 0;
}

void lerPerfil( PerfilSaude *p )
{
    printf( "Nome: " );
    scanf( "%29s", p->nome );

    printf( "Sobrenome: " );
    scanf( "%29s", p->sobrenome );

    printf( "Sexo (M/F): " );
    scanf( " %c", &p->sexo );

    printf( "Data de nascimento (dia mes ano): " );
    scanf( "%d%d%d", &p->nascimento.dia,
                     &p->nascimento.mes,
                     &p->nascimento.ano );

    printf( "Altura em cm: " );
    scanf( "%lf", &p->altura_cm );

    printf( "Peso em kg: " );
    scanf( "%lf", &p->peso_kg );
}

int calcularIdade( const Data *nasc )
{
    time_t t = time( NULL );
    struct tm hoje = *localtime( &t );

    int anoAtual = hoje.tm_year + 1900;
    int mesAtual = hoje.tm_mon + 1;
    int diaAtual = hoje.tm_mday;

    int idade = anoAtual - nasc->ano;

    /* Se ainda nao fez aniversario, subtrai 1 */
    if ( mesAtual < nasc->mes ||
       ( mesAtual == nasc->mes && diaAtual < nasc->dia ) ) {
        idade--;
    }
    return idade;
}

int fcMaxima( int idade )
{
    return 220 - idade;
}

double fcIdealMin( int idade )
{
    return fcMaxima( idade ) * 0.50;
}

double fcIdealMax( int idade )
{
    return fcMaxima( idade ) * 0.85;
}

double calcularIMC( double peso_kg, double altura_cm )
{
    double altura_m = altura_cm / 100.0;
    return peso_kg / ( altura_m * altura_m );
}

const char* classificarIMC( double imc )
{
    if ( imc < 18.5 )       return "Abaixo do peso";
    if ( imc < 25.0 )       return "Normal";
    if ( imc < 30.0 )       return "Acima do peso (sobrepeso)";
    return "Obeso";
}

void exibirPerfil( const PerfilSaude *p )
{
    int idade = calcularIdade( &p->nascimento );
    double imc = calcularIMC( p->peso_kg, p->altura_cm );

    printf( "\n========== PERFIL ==========\n" );
    printf( "Nome:           %s %s\n", p->nome, p->sobrenome );
    printf( "Sexo:           %c\n", p->sexo );
    printf( "Nascimento:     %02d/%02d/%d\n",
            p->nascimento.dia, p->nascimento.mes, p->nascimento.ano );
    printf( "Idade:          %d anos\n", idade );
    printf( "Altura:         %.1f cm\n", p->altura_cm );
    printf( "Peso:           %.1f kg\n", p->peso_kg );

    printf( "\n=== METRICAS DE SAUDE ===\n" );
    printf( "IMC:                  %.2f (%s)\n",
            imc, classificarIMC( imc ) );
    printf( "FC maxima:            %d bpm\n",
            fcMaxima( idade ) );
    printf( "FC ideal (50%%-85%%):  %.0f - %.0f bpm\n",
            fcIdealMin( idade ), fcIdealMax( idade ) );

    printf( "\n=== TABELA DE IMC (OMS) ===\n" );
    printf( "  Abaixo do peso:  IMC < 18.5\n" );
    printf( "  Normal:          18.5 a 24.9\n" );
    printf( "  Sobrepeso:       25.0 a 29.9\n" );
    printf( "  Obeso:           IMC >= 30.0\n" );

    printf( "\nLembre: estes valores sao orientativos.\n" );
    printf( "Sempre consulte um profissional de saude.\n" );
}
\end{lstlisting}

\textbf{Exemplo de execução:}
\begin{saida}
===== SISTEMA DE PERFIL DE SAUDE =====\\
\\
Nome: Maria\\
Sobrenome: Silva\\
Sexo (M/F): F\\
Data de nascimento (dia mes ano): 15 6 1990\\
Altura em cm: 165\\
Peso em kg: 60\\
\\
========== PERFIL ==========\\
Nome:           Maria Silva\\
Sexo:           F\\
Nascimento:     15/06/1990\\
Idade:          35 anos\\
Altura:         165.0 cm\\
Peso:           60.0 kg\\
\\
=== METRICAS DE SAUDE ===\\
IMC:                  22.04 (Normal)\\
FC maxima:            185 bpm\\
FC ideal (50\%-85\%):  93 - 157 bpm\\
\\
=== TABELA DE IMC (OMS) ===\\
\ \ Abaixo do peso:  IMC < 18.5\\
\ \ Normal:          18.5 a 24.9\\
\ \ Sobrepeso:       25.0 a 29.9\\
\ \ Obeso:           IMC >= 30.0
\end{saida}
\end{respostabox}

\subsection*{Análise didática}

\begin{boaprática}
\textbf{Por que \texttt{struct} brilha aqui?}

Sem \texttt{struct}, precisaríamos passar 8 ou mais parâmetros para cada função:
\begin{lstlisting}
/* SEM struct -- horrivel */
void exibir( char *nome, char *sobrenome, char sexo,
             int dia, int mes, int ano,
             double altura, double peso );
\end{lstlisting}

Com \texttt{struct}, encapsulamos tudo:
\begin{lstlisting}
/* COM struct -- limpo */
void exibir( const PerfilSaude *p );
\end{lstlisting}

\textbf{Outros benefícios:}
\begin{itemize}[noitemsep]
  \item Adicionar campos (ex.: alergias, medicamentos) não muda assinaturas de funções.
  \item Passa-se um único ponteiro -- eficiente para estruturas grandes.
  \item Código mais legível e auto-documentado.
\end{itemize}
\end{boaprática}

\bigskip

\begin{atencao}
\textbf{Aninhamento de \texttt{struct}:} note como \texttt{PerfilSaude} contém
\texttt{Data nascimento}. Isso é \textit{composição} -- estruturas dentro de
estruturas. Para acessar campos aninhados via ponteiro:
\begin{lstlisting}
p->nascimento.ano    /* p eh ponteiro; nascimento eh subestrutura */
\end{lstlisting}
A intuição: \texttt{->} ``deszipa'' um ponteiro; \texttt{.} acessa membros
de uma estrutura já desreferenciada.
\end{atencao}

\bigskip

\subsection*{Por que isso é ``fazer a diferença''?}

\begin{tcolorbox}[colback=ecolightgreen,colframe=ecogreen,
  title={\bfseries Registros eletrônicos de saúde}]

A informatização de registros médicos é uma das maiores transformações da
medicina moderna:

\textbf{Benefícios:}
\begin{itemize}[noitemsep]
  \item \textbf{Salva vidas em emergências:} médicos acessam histórico completo
  do paciente -- alergias, medicamentos, condições prévias.
  \item \textbf{Evita interações medicamentosas:} sistemas podem alertar
  automaticamente sobre combinações perigosas.
  \item \textbf{Reduz custos:} elimina exames duplicados, agilizando o atendimento.
  \item \textbf{Pesquisa epidemiológica:} dados agregados (anonimizados) revelam
  padrões de doenças em populações.
\end{itemize}

\textbf{Desafios:}
\begin{itemize}[noitemsep]
  \item \textbf{Privacidade:} dados de saúde são extremamente sensíveis. LGPD
  (Brasil) e HIPAA (EUA) regulam estritamente seu uso.
  \item \textbf{Segurança:} sistemas hospitalares são alvos preferenciais de
  ataques cibernéticos.
  \item \textbf{Interoperabilidade:} sistemas diferentes precisam se comunicar
  (padrões como HL7 FHIR).
  \item \textbf{Inclusão digital:} nem todos têm acesso a tecnologia.
\end{itemize}

Você acabou de construir, em escala minimalista, o núcleo de um sistema
\textbf{Electronic Health Record (EHR)}. Empresas como Epic, Cerner, e a brasileira
MV Sistemas movimentam bilhões de reais com esse tipo de software.
\end{tcolorbox}

\bigskip

% ============================================================
\section*{Reflexão Final -- Capítulo 10}

\begin{tcolorbox}[colback=deitellightblue,colframe=deitelblue,
  title={\bfseries Da memória bruta aos tipos próprios}]

O Capítulo~10 representa um salto qualitativo em sua jornada com C:

\textbf{O que mudou?}
\begin{itemize}[noitemsep]
  \item Antes: variáveis isoladas, arrays homogêneos, ponteiros.
  \item Agora: \textbf{tipos próprios}, modelando entidades do mundo real.
\end{itemize}

\textbf{Conceitos que se tornam possíveis:}
\begin{itemize}[noitemsep]
  \item \textbf{Modelagem de domínio:} \texttt{PerfilSaude}, \texttt{Card},
  \texttt{Conta}, \texttt{Aluno}\ldots
  \item \textbf{Encapsulamento:} função recebe uma estrutura, não 10 parâmetros.
  \item \textbf{Composição:} \texttt{struct} dentro de \texttt{struct} (ex.: \texttt{Data} em \texttt{Perfil}).
  \item \textbf{Manipulação de bits:} controle fino sobre representação binária --
  essencial em sistemas embarcados, gráficos, criptografia.
  \item \textbf{Uniões:} mesmo espaço, múltiplas interpretações -- útil em parsers
  e drivers.
\end{itemize}

\textbf{Próximos passos:}
\begin{itemize}[noitemsep]
  \item \textbf{Cap.~11:} arquivos -- persistir nossas \texttt{struct}s em disco.
  \item \textbf{Cap.~12:} ponteiros para \texttt{struct} formam listas encadeadas,
  filas, pilhas, árvores -- as estruturas de dados clássicas.
  \item \textbf{Cap.~13:} pré-processador para macros poderosas com estruturas.
\end{itemize}

A combinação \texttt{struct + typedef + ponteiro} é o tripé sobre o qual se
constrói todo o software sério em C. Domine esses três e você terá as ferramentas
para projetar sistemas de qualquer complexidade.

\end{tcolorbox}

\end{document}
